\documentclass[11pt]{article}
\usepackage[margin=1in]{geometry}
\usepackage{booktabs}
\usepackage{graphicx}
\usepackage{hyperref}
\usepackage{amsmath}
\usepackage{natbib}
\usepackage{xcolor}

\title{Length-Dependent Accuracy and Construct Mismatch in LLM-as-Judge:\\Lessons from Step-Level Agent Degradation Measurement}
\author{Alex Reff\\Independent Researcher\\\texttt{alex.reff@gmail.com}}
\date{}

\begin{document}
\maketitle

\begin{abstract}
A measurement pipeline grades each step of a coding-agent trace with an LLM judge, validates the judge against human labels, decomposes the degradation slope from confounds, and ships an optional SIMEX label-noise correction calibrated to the grader's TRAIL-validated flip rate.
Across 14 BH-corrected configurations (one further configuration was not re-graded; 8 models, 4 scaffoldings, ${\sim}$24,000 graded steps total, full grader context, parse-error steps excluded), 5 show statistically significant degradation (all 5 surviving Benjamini--Hochberg correction at $\alpha=0.05$), 2 show improvement (1 surviving BH), and 7 are null.
The dominant confound is phase composition---agents shift from exploration to action over a task, and that shift mimics degradation---but it attenuates rather than eliminates the signal on configurations with larger samples.
A long-trace follow-up (50 traces, 40+ steps, three Llama variants), designed specifically to test context-window pressure, shows $+0.0007$ errors/step ($p < 0.001$, BH-significant).
Two findings about the grading instrument itself generalize beyond this study.
First, grader $\kappa$ against human labels decays with cumulative input length rather than step position: it peaks at 10--50K characters and falls past 50K for both MiniMax and Haiku, consistent with general long-context accuracy degradation.
Second, the 73\% false-negative rate at MEDIUM+ severity is consistent with a construct mismatch---the grader evaluates decision quality at the time of action while human reference labels effectively encode outcome contribution---rather than grader failure; $\kappa$ reaches 0.49 at the HIGH-only threshold where errors are recognizable without future context. The size of the mismatch is bounded by the kappa gap; a registered hindsight-vs-blind two-pass experiment is pending direct measurement.
The pipeline and all graded traces are released as open-source tools.
\end{abstract}

\section{Introduction}

LLM-as-judge is increasingly used to evaluate agent behavior at the step level---grading each action in a trace for correctness, relevance, or progress.
But if the judge's accuracy varies systematically with position in the trace, any temporal analysis built on those judgments inherits a bias the researcher cannot see.
This paper documents such an effect and traces its consequences through a study of step-level degradation in coding agents.

Agent benchmarks---SWE-bench \citep{jimenez2024swebench}, GAIA \citep{mialon2023gaia}, WebArena \citep{zhou2024webarena}---measure final outcomes: did the agent solve the task?
They do not measure what happens \emph{during} a task.
An agent that makes one early mistake and flawlessly executes the wrong plan for 20 steps looks identical to one that independently fails every step.
These are different failure modes requiring different interventions, and no existing tooling distinguishes them.

The assumption that agents degrade over long tasks appears widely in framework documentation and engineering practice.
Teams build context management into agent scaffolding---sliding windows that drop early messages, summarization chains that compress history, periodic agent restarts that discard state---partly justified by the assumption that long contexts degrade accuracy.
Long-context degradation is well-documented for \emph{retrieval} tasks \citep{liu2024lost}, but retrieving a planted fact from a long prompt is a different capability than generating a correct next action given a history of prior actions.
There does not appear to be published work measuring within-run step-level error trajectories for coding agents specifically.

This paper presents a study that measured step-level degradation across 15 agent configurations, identified the dominant confound (phase composition), and characterized two properties of the grading instrument---length-dependent accuracy and a construct mismatch with human reference labels---that affect any temporal analysis built on LLM-as-judge outputs.
The contributions are:

\begin{enumerate}
    \item A 14-configuration degradation study (${\sim}$24,000 graded steps): of the 5 with statistically significant degradation, all 5 survive BH-FDR correction at $\alpha=0.05$; 2 configurations show improvement (1 surviving BH); 7 are null. Degradation is not model-specific---it appears across GPT-4o, Claude 3.5 Sonnet, Claude 3.7 Sonnet, and Llama (predominantly 70B in the long-trace sample). Phase composition is a real confound that attenuates slopes by roughly half but does not eliminate them on configurations with larger samples or longer traces.
    \item Evidence that grader $\kappa$ decays with \emph{input length} rather than step position: stratifying TRAIL validation by cumulative prior-step character count, $\kappa$ peaks at 10--50K chars and decays past 50K for both MiniMax and Haiku, consistent with general long-context accuracy degradation rather than intrinsic position bias.
    \item A construct mismatch between no-hindsight grading and hindsight-informed reference labels: $\kappa$ against TRAIL's expert annotations systematically underestimates grader accuracy on the decision-quality construct the grader actually measures. The 73\% false-negative rate at MEDIUM+ likely reflects errors identifiable only with future context.
    \item A reusable measurement pipeline (\texttt{inspect-degradation}) integrating step-level LLM grading, grader validation against human labels, optional SIMEX label-noise correction, and confound decomposition, released as an open-source Inspect AI extension.
\end{enumerate}

\section{Related Work}

\textbf{Agent benchmarks and reliability.}
SWE-bench \citep{jimenez2024swebench}, GAIA \citep{mialon2023gaia}, and WebArena \citep{zhou2024webarena} evaluate task-level success.
The HAL leaderboard \citep{kapoor2026hal} provides standardized multi-benchmark evaluation with cost tracking, and its Reliability Dashboard measures consistency, robustness, and predictability across runs---but not within-run temporal dynamics.
The CLEAR framework \citep{mehta2025clear} proposes multi-dimensional enterprise evaluation across cost, latency, efficacy, assurance, and reliability.
None of these measure within-run temporal dynamics.
TRAIL \citep{patronus2024trail} provides expert-annotated step-level labels for agent traces; this work extends TRAIL's annotations to temporal trajectory analysis by fitting mixed-effects models to step-level grades across the trace timeline.
Anthropic's ``Demystifying Evals'' \citep{anthropic2025evals} discusses per-step evaluation methodology; this work provides a reusable implementation and reports the degradation measurements that methodology enables.

\textbf{LLM-as-judge.}
MT-Bench and Chatbot Arena \citep{zheng2023judging} evaluate single responses or conversations; AlpacaEval \citep{alpacaeval} evaluates instruction-following quality.
This work extends the LLM-as-judge paradigm to multi-step agent traces where each grading call receives a growing context of prior steps, and adds measurement-error correction (SIMEX) to account for the noise the judge introduces.
Kapoor et al.\ \citep{kapoor2024scalable} found that cheap-judge confidence is unreliable as an escalation gate in cascade architectures, motivating this work's use of ensemble voting \citep{wang2022selfconsistency} over cascade designs and the exclusion of verbalized confidence scores from the rubric.

\textbf{Long-context retrieval.}
Liu et al.\ \citep{liu2024lost} showed that LLMs struggle to use information placed in the middle of long contexts (``Lost in the Middle'').
Needle-in-a-haystack benchmarks measure fact retrieval as context grows.
These are real effects, but they test retrieval from context, not action generation given context.

\textbf{Measurement error correction.}
SIMEX \citep{cook1994simulation} corrects regression coefficients for covariate measurement error.
It is adapted here for binary label noise in LLM-graded outcomes, following the framework of Carroll et al.\ \citep{carroll2006measurement}.

\section{Method}

\subsection{Step-Level Grading}

Each step of an agent trace is graded by an LLM along five dimensions using a structured rubric:

\begin{itemize}
    \item \textbf{Validity} (fail / neutral / pass): Does the step advance the task? ``Pass'' requires positive, pointable evidence of progress---not merely the absence of visible error.
    \item \textbf{Complexity} (low / medium / high): How clear was the correct next move? A confound control variable.
    \item \textbf{Dependency} (independent / dependent / n/a): Is a failure caused by a prior failure? Enables cascade-chain analysis.
    \item \textbf{Severity} (low / medium / high): Recovery cost of a failing step.
    \item \textbf{Looping} (boolean): Is this step a repeat of a prior action? Mechanical, occurrence-based.
\end{itemize}

All dimensions are closed-set categoricals.
No confidence scores are elicited---verbalized LLM self-confidence is unreliable \citep{xiong2023can}.
Uncertainty is measured empirically via validation, not self-reported.

The grader prompt includes all prior steps in the trace as context.
The pipeline supports a configurable \texttt{prior\_context\_char\_budget} that truncates prior-step content.
An initial run applied a 30,000-character cap for cost control; this was subsequently identified as introducing position-correlated bias (truncation rate scaled from 0\% at step 0 to 88\% at step 25), and all 14 affected configurations were re-graded with full context.
All results in this paper use uncapped grading with parse-error steps excluded.
See Appendix~\ref{app:correction} for the correction details.
Parse failures (truncated JSON, rate-limit exhaustion, provider errors) fall back to \texttt{Validity.neutral} in the library; these steps are excluded from all analyses reported here.

\subsection{Step-Phase Classification}

A deterministic classifier labels each step as \emph{exploration} (reads, searches, navigation) or \emph{action} (edits, test runs, submissions) from the agent's command text.
Framework-specific detection layers handle SWE-agent XML blocks, OpenHands bracket commands with subcommand parsing, terminus XML, Auto-SWE structured tool calls, and a shell-command fallback.
Validation: 100\% accuracy on Auto-SWE (verified against ground-truth tool names), 98--100\% on SWE-agent (100-step human-reviewed sample).

\subsection{Statistical Analysis}

The step-level grades feed a mixed-effects linear probability model:
\[
\text{is\_error}_{ij} = \beta_0 + \beta_1 \cdot \text{step\_index}_{ij} + \beta_2 \cdot \text{phase}_{ij} + \beta_3 \cdot \text{complexity}_{ij} + \beta_4 \cdot \text{outcome}_i + u_i + \varepsilon_{ij}
\]
where $i$ indexes traces (with random intercept $u_i$ grouped by task), $j$ indexes steps, and $\beta_1$ is the degradation coefficient of interest.
The model adapts to available covariates: outcome is dropped when unavailable, model identity is added when multiple models are present.
Confidence intervals on rates, slopes, and chain-length statistics use trace-level bootstrap resampling (whole traces resampled, not steps) to correctly account for within-trace correlation. Mixed-effects regression coefficients use the asymptotic Wald CI from the fitted random-intercept model, where the random intercept absorbs the trace-level correlation directly.
Benjamini--Hochberg FDR correction is applied across the configuration family.

\subsection{Grader Selection and Validation}
\label{sec:method:validation}

Seven grader configurations were evaluated against TRAIL's 148 expert-annotated traces (954 step pairs).
Cohen's $\kappa$ was computed at three severity thresholds: any error, MEDIUM+, and HIGH-only.
Four of five graders reached $\kappa \approx 0.45$--$0.50$ at the HIGH-only threshold (moderate agreement), with disagreement concentrating on LOW-impact errors.
Three additional findings emerged from validation:

\emph{Ensembles do not improve agreement.}
Neither 3-sample self-consistency nor cross-family majority voting (Haiku + MiniMax + Gemini) beat the best single model at any threshold.

\emph{Frontier models do not improve grading.}
Sonnet 4.6, Kimi K2.5, and Gemini 2.5 Flash Lite all underperformed MiniMax and Haiku.

\emph{Rubric iteration has negative returns.}
Adding a cosmetic-error exception collapsed binary $\kappa$ from 0.202 to 0.029.
The exception language gave the grader an escape hatch it used far beyond its intended scope.

MiniMax M2.5 (\$0.40/M input tokens) was selected as the primary grader: $\kappa = 0.49$ at HIGH-only, matching Haiku-level accuracy at lower cost.
The SIMEX flip probability is derived from the HIGH-only confusion matrix (flip rate = 0.12), which matches the grader's effective operating threshold. The correction is available in the released pipeline; it is not applied to the headline coefficients in this paper because at this flip rate the implied slope attenuation is small relative to the cross-configuration variance, and the registered Lock D follow-up will produce the corpus needed to choose between flip-rate calibrations.

\section{Experimental Setup}

Fifteen configurations spanning 8 distinct models and 4 scaffoldings were tested, totaling ${\sim}$24,000 graded steps (Table~\ref{tab:configs}).
Most use SWE-bench Python bug-fix instances; one (Auto-SWE) uses a custom multi-agent pipeline on production tasks.
Several models appear in multiple configurations: ``Crossover'' runs test a model on a scaffolding different from its original dataset, while ``MSB'' runs use Multi-SWE-bench trajectories that provide the same model--scaffolding pair on a different task sample.
This design separates model effects from scaffolding effects and tests replication across task distributions.

All coefficients in Table~\ref{tab:configs} are from \textbf{uncapped} grader caches with parse-error steps excluded, following the methodology correction described in Appendix~\ref{app:correction}.
One configuration (MSB / GPT-4o / OpenHands) had 0\% truncation under the original cap and was not re-graded.

\begin{table}[h]
\centering
\small
\caption{Cross-dataset summary (post-correction). All coefficients from mixed-effects LPM on uncapped grader caches, parse-error steps excluded, random intercept for task, controlling for step phase, complexity, and outcome where available. Five slopes survive BH-FDR correction across the 14-pair family.}
\label{tab:configs}
\begin{tabular}{llrrl}
\toprule
Model & Scaffolding & Traces & $\beta_1$ (step\_index) & Direction \\
\midrule
Claude 3.5 Sonnet & SWE-agent (Crossover) & 50 & $+0.0016^{*}$ & Degrades \\
Claude 3.5 Sonnet & SWE-agent (MSB) & 50 & $+0.0010$ & No effect \\
Claude 3.5 Sonnet & OpenHands (Crossover) & 50 & $+0.0007$ & No effect \\
Claude 3.5 Sonnet & OpenHands (MSB) & 50 & $+0.0029^{***}$ & Degrades \\
Claude 3.7 Sonnet & SWE-agent & 30 & $+0.0003$ & No effect \\
Claude 3.7 Sonnet & OpenHands (MSB) & 50 & $+0.0014^{**}$ & Degrades \\
GLM 4.7 & terminus-2 & 50 & $-0.0059^{***}$ & Improves \\
GPT-4o & SWE-agent (Crossover) & 50 & $+0.0003$ & No effect \\
GPT-4o & SWE-agent (MSB) & 50 & $+0.0025^{**}$ & Degrades \\
GPT-4o & OpenHands & 50 & $-0.0013$ & No effect \\
GPT-4o & OpenHands (MSB)$^\dagger$ & 50 & $-0.0003^{*}$ & Improves \\
Llama 70B & SWE-agent & 30 & $+0.0030$ & No effect \\
Llama (all, 40+ steps) & SWE-agent & 50 & $+0.0007^{***}$ & Degrades \\
Qwen3-Coder-480B & OpenHands & 30 & $+0.0002$ & No effect \\
Qwen3-Coder-Next & Custom pipeline & 50 & $-0.0009^{*}$ & Improves \\
\bottomrule
\multicolumn{5}{l}{\footnotesize $^{*}p < 0.05$, $^{**}p < 0.01$, $^{***}p < 0.001$. $^\dagger$Not re-graded (0\% truncation).}
\end{tabular}
\end{table}

\section{Results}

\subsection{Phase-Composition Confound}

The first dataset analyzed (Nebius / Llama 70B, 30 traces, 632 steps) showed apparent degradation.
Each successive control attenuated the signal.
Under the corrected methodology (uncapped grading, parse-error steps excluded), the phase-controlled slope sits at $+0.0030$/step ($p = 0.055$, 95\% CI $[-0.00007, +0.00603]$)---borderline null:

\begin{table}[h]
\centering
\caption{Progressive confound control on Nebius / Llama 70B (uncapped).}
\label{tab:dismantling}
\begin{tabular}{lrrl}
\toprule
Controls & $\beta_1$ & $p$-value & What it revealed \\
\midrule
Raw per-trace OLS & ${\sim}+0.020$ & varies & --- \\
+ Mixed-effects & positive & --- & Task variance, outcome bias \\
+ Step-phase covariate & $+0.0030$ & $0.055$ & Phase attenuates but does not eliminate \\
\bottomrule
\end{tabular}
\end{table}

Agents explore early (reads, searches---low error rate) and act late (edits, test runs---high error rate).
Step phase is the dominant predictor of error rate: action steps have 11--30 percentage points higher error rates than exploration steps ($p < 0.0001$ in every configuration), an effect larger than any model or task covariate.
Controlling for phase approximately halves the raw slope.
But it does not drive it to zero---on configurations with larger samples or longer traces (the long-trace follow-up, four re-graded MSB configurations), the phase-controlled slope remains significantly positive.

\subsection{Cross-Dataset Summary}

Five of 14 configurations show statistically significant degradation after phase control (Table~\ref{tab:configs}).
The pattern is concentrated in Multi-SWE-bench and long-trace data: 3 of 4 MSB configurations, the long-trace follow-up, and Crossover/Claude 3.5/SWE-agent.
Two configurations show improvement (Terminus/GLM 4.7, Auto-SWE/Qwen3-Coder-Next).
Seven are null.

Degradation is not model-specific: GPT-4o, Claude 3.5 Sonnet, Claude 3.7 Sonnet, and Llama (predominantly 70B; the Nebius long sample is 41 traces of 70B, 8 of 8B, 1 of 405B) all show degradation on at least one configuration.
Qwen3-Coder-480B is the only model with no degradation signal across any framework tested.
Three configurations (OpenHands/GPT-4o, OpenHands/Qwen3-Coder, SWE-smith/Claude 3.7) are consistently null.

\subsection{Long-Trace Follow-Up}

A targeted follow-up tested whether degradation appears under maximum context-window pressure, using 50 traces with 40+ steps each (3,800 steps) across three Llama variants (70B, 8B, 405B).
Result: $\beta_1 = +0.0007$ ($p < 0.001$, survives BH-FDR).
Over a 100-step trace this adds 7 percentage points in expected error rate---non-negligible relative to the 26\% baseline.
This is the single configuration most directly testing the context-window-pressure hypothesis, and the signal is the strongest in the study.

\subsection{Grader Sensitivity}

The same 30 Nebius traces were re-graded with Claude Haiku 4.5 to test whether the degradation slope is sensitive to grader choice:

\begin{table}[h]
\centering
\caption{Grader sensitivity test on Nebius / Llama 70B.}
\label{tab:sensitivity}
\begin{tabular}{lrr}
\toprule
Grader & $\beta_1$ (with phase) & $p$-value \\
\midrule
MiniMax & $+0.0030$ & $0.055$ \\
Haiku & $+0.0117$ & $< 0.001$ \\
\bottomrule
\end{tabular}
\end{table}

Both graders trend positive.
Haiku finds significant degradation; MiniMax is borderline null (CI crosses zero by 0.00007).
The ${\sim}4\times$ magnitude gap is real---the two models calibrate the ``fail'' boundary differently---but the direction of effect is consistent.
A second grader producing a divergent slope is a cheap diagnostic for grader-specific bias (under \$1 to re-grade 632 steps with Haiku); here the divergence is in magnitude rather than direction, suggesting both graders are measuring the same underlying signal at different sensitivities.

\subsection{Severity-Threshold Calibration}

$\kappa$ increases monotonically with severity threshold across all five grader configurations tested (Table~\ref{tab:severity}).

\begin{table}[h]
\centering
\caption{$\kappa$ by severity threshold.}
\label{tab:severity}
\begin{tabular}{lrrrrr}
\toprule
Threshold & Haiku & MiniMax & MiniMax SC3 & Trio & Gemini \\
\midrule
Any error & 0.221 & 0.202 & 0.212 & 0.215 & 0.061 \\
MEDIUM+ & 0.278 & 0.251 & 0.259 & 0.266 & 0.065 \\
HIGH only & \textbf{0.450} & \textbf{0.486} & \textbf{0.503} & \textbf{0.474} & 0.105 \\
\bottomrule
\end{tabular}
\end{table}

At MEDIUM+, the false-negative rate is 73\% overall and 91\% at later steps.
The grader effectively detects HIGH-severity errors only, regardless of rubric language (Figure~\ref{fig:severity}).

\begin{figure}[h]
\centering
\includegraphics[width=0.7\textwidth]{figures/severity_threshold.png}
\caption{Cohen's $\kappa$ by severity threshold across five grader configurations. Agreement increases monotonically with threshold strictness, reaching moderate levels ($\kappa \approx 0.45$--$0.50$) only at HIGH-only. The pattern is consistent across model families.}
\label{fig:severity}
\end{figure}

\subsection{Construct Mismatch: Decision Quality vs.\ Outcome Contribution}
\label{sec:construct}

The 73\% false-negative rate at MEDIUM+ conflates two constructs.
The grader operates under a no-hindsight constraint: it judges each step using only prior context, without seeing future steps.
TRAIL's human annotators labeled traces after the fact, with full knowledge of consequences.

An edit that compiles and looks correct at step 3 but turns out to be wrong when a test fails at step 60 is a \emph{good decision} (the grader's construct) but a \emph{bad outcome contribution} (TRAIL's construct).
This is not grader error---it is a construct mismatch between decision-quality evaluation and outcome-contribution evaluation.

$\kappa$ against hindsight-informed reference labels systematically underestimates grader accuracy on the construct the grader is actually measuring.
At the HIGH-only threshold---where errors are severe enough to be recognizable at decision time (corrupted state, hallucinated facts, explicitly violated instructions)---$\kappa$ reaches 0.49, consistent with moderate agreement on a genuinely measurable construct.

This distinction generalizes: any LLM-as-judge rubric with a no-hindsight constraint, validated against hindsight-informed labels, will show artificially low agreement.
TRAIL's category labels (e.g., ``Instruction Non-compliance'') are often too underspecified to support formal classification of false negatives by hindsight-dependence, and $\kappa$ against TRAIL may represent a lower bound on inter-rater agreement with any reviewer, including the TRAIL annotators themselves.

\subsection{Length-Dependent Grader Accuracy}

Binary $\kappa$ against TRAIL drops from 0.33 (steps 0--2) to 0.03 (steps 6+).
The original framing attributed this to ``position-dependent grader conservatism.''
Stratifying the same data by cumulative prior-step character count reveals a different mechanism---$\kappa$ decays with input length:

\begin{table}[h]
\centering
\caption{$\kappa$ by prior-step character count (MiniMax, $n=954$).}
\label{tab:length_kappa}
\begin{tabular}{lrr}
\toprule
Prior-step length & $n$ & $\kappa$ \\
\midrule
$<$10K & 344 & $+0.147$ \\
10--50K & 361 & $\mathbf{+0.226}$ \\
50--200K & 238 & $+0.094$ \\
$>$200K & 11 & $+0.029$ \\
\bottomrule
\end{tabular}
\end{table}

Haiku shows the same pattern ($\kappa$: $+0.275$ at 10--50K $\to$ $+0.059$ at 50--200K $\to$ $-0.375$ at $>$200K), consistent with long-context accuracy decay as a general LLM property.

Position and length are heavily confounded in TRAIL: 339/438 steps at position 0--2 have $<$10K prior chars; 154/175 steps at position 10+ have 50--200K.
The diagonal corners of the joint (position $\times$ length) table are effectively empty, so the two cannot be cleanly disentangled at this sample size.
But the length mechanism explains the observed $\kappa$ curve without requiring intrinsic position bias.

\subsection{Power Analysis}

Monte Carlo simulation (80 iterations per cell, flip probability 0.12) on the original capped data suggested the study could detect slopes of ${\sim}0.01$/step at 80\% power for typical corpus sizes (30--50 traces, 15--25 steps).
Under the corrected methodology, several configurations produced statistically significant slopes well below this floor: $+0.0007$ (phase3-long, $p < 0.001$), $+0.0014$ (MSB/Claude 3.7/OpenHands, $p = 0.001$), $+0.0007$ (phase3-autoswe-implement, $p = 0.009$).
The original MDE was conservative---the study's effective sensitivity is higher than the power table implied, likely because removing parse-error fallbacks reduced measurement noise.

\subsection{Improvement Signals}

Under capped grading, 6 of 15 configurations showed significant improvement.
Under uncapped grading, 4 of those 6 retracted (three MSB configurations flipped to degradation; one to null).
Two improvement signals survive the correction: Terminus/GLM 4.7 ($-0.0059$, $p < 0.001$, BH-significant in both regimes) and Auto-SWE/Qwen3-Coder-Next ($-0.0009$, $p = 0.034$, raw-$p$ only).
Neither can be outcome-controlled---Terminus has \texttt{success=None}, and Auto-SWE has a process-completion proxy but not ground-truth resolution---so the improvement is indistinguishable from survivorship selection with the currently available data.

\section{Practical Guidance}

\textbf{Provide the grader full context, or instrument the truncation.}
If capping is required for cost, the cap should be enforced uniformly across positions (e.g., a rolling window that truncates head content rather than tail) and truncation exposure should be modeled as a covariate in downstream regressions.
Position-stratifying accuracy without stratifying by truncation exposure produces misleading conclusions (Appendix~\ref{app:correction}).

\textbf{Handle parse-error steps explicitly.}
Grader libraries that fall back to a default label on parse failure introduce position-correlated bias when the failure rate varies with input length.
Analyses should drop parse-error steps by default and treat their inclusion as an explicit decision.

\textbf{Use a second grader as a diagnostic.}
Running a second grader on the same traces is a cheap diagnostic for grader-specific bias (under \$1 to re-grade 632 steps with Haiku), as long as both graders see matched context.

\textbf{Apply multiple-comparisons correction.}
Fourteen configuration-level regressions in the BH-FDR family at raw $\alpha = 0.05$ produce ${\sim}0.7$ false positives by chance.
This paper applies BH-FDR correction across the configuration family.

\textbf{Stratify grader validation by input length.}
The length-dependent $\kappa$ decay documented in Table~\ref{tab:length_kappa} means that any LLM-as-judge validation should stratify by input length, not just step position.

\section{Limitations}

\begin{itemize}
    \item \textbf{Measurement instrument quality.} The primary grader achieves $\kappa = 0.49$ at HIGH-only---moderate agreement on errors detectable at decision time. The 73\% FNR at MEDIUM+ likely reflects hindsight-dependent labels in TRAIL rather than grader failures (Section~\ref{sec:construct}), but this cannot be confirmed without decomposing TRAIL's labels by hindsight-dependence. TRAIL's inter-annotator agreement is unpublished, so $\kappa$ values lack an interpretive ceiling. The rubric has not been validated by human experts independent of the TRAIL labels.
    \item \textbf{Context management unknown.} SWE-agent and OpenHands have configurable context management, and the trajectory data does not record which settings were used. If a framework silently drops context, step\_index no longer reflects context length.
    \item \textbf{Step granularity is not comparable across datasets.} OpenHands steps are multi-tool-call groups; Nebius and Auto-SWE steps are single LLM calls. Cross-dataset claims about \texttt{step\_index} robustness are comparing incomparable axes.
    \item \textbf{Narrow task distribution.} Mostly SWE-bench Python bug fixes. 30--50 traces per configuration.
    \item \textbf{Methodological correction.} An initial run used a 30K-character context cap that introduced position-correlated bias; all results reported here use uncapped grading (Appendix~\ref{app:correction}). Some subsidiary analyses (cascade chain length, autocorrelation, phase-stratified regressions) have not been re-estimated against the uncapped data.
\end{itemize}

\section{Conclusion}

Naive step-level regressions on agent traces produce convincing degradation artifacts.
The dominant confound is phase composition: agents shift from exploration to action over a task, and action steps have higher error rates.
But phase composition attenuates rather than eliminates the degradation signal: 5 of 14 configurations show significant degradation after phase control, including the long-trace follow-up ($+0.0007$/step, $p < 0.001$, BH-significant) designed specifically to test context-window pressure.

Two findings about the grading instrument generalize beyond this study.
First, grader $\kappa$ decays with input length rather than step position---peaking at 10--50K characters and falling past 50K---which means any temporal analysis using LLM judges on traces with growing context inherits a length-correlated accuracy gradient.
Second, the 73\% false-negative rate at MEDIUM+ severity reflects a construct mismatch between decision-quality evaluation (the grader's construct) and outcome-contribution evaluation (TRAIL's construct), not grader failure.
At the HIGH-only threshold, where errors are recognizable without future context, $\kappa$ reaches 0.49.
This distinction applies to any LLM-as-judge rubric with a no-hindsight constraint validated against hindsight-informed labels.

\textbf{Implications.}
For teams building agent evaluation systems: LLM-as-judge pipelines applied to multi-step traces should provide the grader full context (or instrument any truncation as a downstream covariate), exclude parse-error steps by default, stratify grader validation by input length, and apply multiple-comparisons correction across configuration families.
For teams building context management into agent scaffolding: the evidence that agents degrade over long tasks is not zero---but 7 of 14 configurations are null, and where present the degradation magnitude is typically 0.001--0.003 errors/step.
Context-management decisions should be evaluated on both cost/latency benefits and on whether the downstream evaluation can distinguish real temporal signals from measurement artifacts.

The measurement pipeline (\texttt{inspect-degradation}) and all graded traces are available at \url{https://github.com/reffdev/inspect-degradation} and \url{https://github.com/reffdev/inspect-degradation-study}.

\begin{ack}
This work was conducted independently without institutional affiliation or external funding. All compute costs were self-funded. The author thanks the teams behind TRAIL, SWE-bench, Multi-SWE-bench, and the Nebius trajectory release for making their data publicly available.
\end{ack}

\appendix

\section{Grader Rubric (v1)}
\label{app:rubric}

The complete system prompt and user template used for all grading runs are reproduced below.
The rubric is stored as YAML in the \texttt{inspect-degradation} package and loaded at runtime; the version field enables correlation of cached grades with the rubric that produced them.

\subsection*{System Prompt}

\begin{quote}
\small
\begin{verbatim}
You are a step-level evaluator for AI agent traces. You grade ONE
step - the CURRENT step - against the original task goal, using
the prior steps as context. Your output is one JSON object,
nothing else.

## The rule that matters most

The hardest failure mode of this rubric is over-crediting a stuck
agent. "pass" does NOT mean "I see no problem". It means:

    "I can complete this sentence with a concrete noun phrase:
     'This step advanced the task by producing ________.'  "

If the best you can say is "the agent thought about the problem"
or "stated its plan" or "made a tool call" without naming the
concrete fact, decision, or artifact that came out of it - that
is NOT a pass. It is "neutral".

Default toward "neutral" whenever you would otherwise reach for
"pass" on autopilot. The rubric is deliberately biased this way,
because under-crediting is recoverable and over-crediting is not.

## Fields

validity - does this step visibly advance the task?
  - "fail"    : you can point to a specific thing that is wrong.
  - "neutral" : not wrong, but you cannot point to concrete
                progress. Default for any step where you would
                reach for "pass" on autopilot.
  - "pass"    : visible, pointable progress.

complexity - how clear was the correct next move BEFORE this step?
  - "low" / "medium" / "high"

dependency - only meaningful when validity=="fail".
  - "independent" / "dependent" / "n/a"

severity - only meaningful when validity=="fail".
  Anchored to recovery cost: "low" / "medium" / "high"

is_looping - is this step at least the SECOND occurrence of a
  substantially identical action within the prior steps?
  - true / false

## Output

Return exactly one JSON object with these keys:
    validity, complexity, dependency, severity, is_looping
\end{verbatim}
\end{quote}

\noindent
The full rubric includes additional detail on cross-step signals (repetition, self-contradiction, fabricated verification, hedging loops, phantom progress), hard rules (no hindsight, judge the step not the agent, tie-break downward), concrete categories for each validity level, and a pre-response checklist in the user template.
The complete YAML is available in the repository at \texttt{src/inspect\_degradation/prompts/step\_grader\_v1.yaml}.

\section{Methodology Correction: Context-Cap Removal}
\label{app:correction}

An initial run of the grading pipeline applied \texttt{prior\_context\_char\_budget=30000} to every grader call for cost control.
Truncation rate was position-correlated: 0\% at step 0, scaling to 88\% at step 25 in the worst configuration, with four configurations exceeding 60\% average truncation.
Truncated completions produced parse errors that the grader library converted to \texttt{Validity.neutral} fallbacks, concentrating neutral labels at later positions and suppressing late-step error detection.

All 14 affected configurations were re-graded with \texttt{prior\_context\_char\_budget=None} and re-analyzed with parse-error steps excluded.
Six configurations changed direction of effect or significance under the correction:

\begin{table}[h]
\centering
\small
\caption{Effect of removing the 30K context cap.}
\label{tab:correction}
\begin{tabular}{lll}
\toprule
Configuration & Before correction & After correction \\
\midrule
Long-trace (Llama, 40+ steps) & No effect, $+0.0001$ $p{=}0.38$ & Degrades, $+0.0007$ $p{<}0.001$ \\
MSB / GPT-4o / SWE-agent & Improves, $p{<}0.0001$ & Degrades, $+0.0025$ $p{=}0.001$ \\
MSB / Claude 3.5 / OpenHands & Improves, $p{<}0.0001$ & Degrades, $+0.0029$ $p{<}0.001$ \\
MSB / Claude 3.7 / OpenHands & Improves, $p{=}0.007$ & Degrades, $+0.0014$ $p{=}0.001$ \\
MSB / Claude 3.5 / SWE-agent & Improves, $p{<}0.0001$ & No effect, $+0.0010$ $p{=}0.12$ \\
Auto-SWE / Qwen3-Coder-Next & No effect, $p{=}0.11$ & Improves, $-0.0009$ $p{=}0.034$ \\
\bottomrule
\end{tabular}
\end{table}

Three configurations were null under both regimes; two retained their original sign and significance.
All results in the main text use the corrected (uncapped) data.
Full cap-vs-uncap comparison data is available at \texttt{results/compare\_all\_pairs\_final.json} in the study repository.

\bibliographystyle{plainnat}
\begin{thebibliography}{99}

\bibitem[Anthropic(2025)]{anthropic2025evals}
Anthropic.
\newblock Demystifying evals for AI agents, 2025.
\newblock URL: \url{https://www.anthropic.com/engineering/demystifying-evals-for-ai-agents}.

\bibitem[Carroll et~al.(2006)]{carroll2006measurement}
R.~J. Carroll, D.~Ruppert, L.~A. Stefanski, and C.~M. Crainiceanu.
\newblock \emph{Measurement Error in Nonlinear Models: A Modern Perspective}.
\newblock Chapman and Hall/CRC, 2nd edition, 2006.

\bibitem[Cook and Stefanski(1994)]{cook1994simulation}
J.~R. Cook and L.~A. Stefanski.
\newblock Simulation-extrapolation estimation in parametric measurement error models.
\newblock \emph{Journal of the American Statistical Association}, 89(428):1314--1328, 1994.

\bibitem[Jimenez et~al.(2024)]{jimenez2024swebench}
C.~E. Jimenez, J.~Yang, A.~Wettig, S.~Yao, K.~Pei, O.~Press, and K.~Narasimhan.
\newblock {SWE}-bench: Can language models resolve real-world {GitHub} issues?
\newblock In \emph{ICLR}, 2024.

\bibitem[Kapoor et~al.(2024)]{kapoor2024scalable}
A.~Kapoor, N.~Raman, Y.~Peng, and S.~Narayanan.
\newblock On scalable oversight with weak {LLMs} judging strong models.
\newblock \emph{arXiv preprint arXiv:2407.04622}, 2024.

\bibitem[Kapoor et~al.(2026)]{kapoor2026hal}
S.~Kapoor et~al.
\newblock {HAL}: Holistic Agent Leaderboard --- the missing infrastructure for AI agent evaluation.
\newblock In \emph{ICLR}, 2026.
\newblock arXiv:2510.11977.

\bibitem[Li et~al.(2024)]{alpacaeval}
X.~Li, T.~Zhang, Y.~Dubois, R.~Taori, I.~Gulrajani, C.~Guestrin, P.~Liang, and T.~B. Hashimoto.
\newblock {AlpacaEval}: An automatic evaluator for instruction-following language models, 2024.
\newblock URL: \url{https://github.com/tatsu-lab/alpaca_eval}.

\bibitem[Liu et~al.(2024)]{liu2024lost}
N.~F. Liu, K.~Lin, J.~Hewitt, A.~Paranjape, M.~Bevilacqua, F.~Petroni, and P.~Liang.
\newblock Lost in the middle: How language models use long contexts.
\newblock \emph{TACL}, 2024.

\bibitem[Mehta et~al.(2025)]{mehta2025clear}
S.~Mehta et~al.
\newblock Beyond accuracy: a multi-dimensional framework for evaluating enterprise agentic AI systems ({CLEAR}).
\newblock arXiv:2511.14136, 2025.

\bibitem[Mialon et~al.(2023)]{mialon2023gaia}
G.~Mialon et~al.
\newblock {GAIA}: A benchmark for general {AI} assistants.
\newblock \emph{arXiv preprint arXiv:2311.12983}, 2023.

\bibitem[Patronus AI(2024)]{patronus2024trail}
Patronus~AI.
\newblock {TRAIL}: Trace reasoning and agentic issue localization, 2024.
\newblock URL: \url{https://github.com/patronus-ai/trail-benchmark}.

\bibitem[Wang et~al.(2023)]{wang2022selfconsistency}
X.~Wang, J.~Wei, D.~Schuurmans, Q.~Le, E.~Chi, S.~Narang, A.~Chowdhery, and D.~Zhou.
\newblock Self-consistency improves chain of thought reasoning in language models.
\newblock In \emph{ICLR}, 2023.

\bibitem[Xiong et~al.(2023)]{xiong2023can}
K.~Xiong, X.~Ding, Y.~Cao, T.~Liu, and B.~Qin.
\newblock Can {LLMs} express their uncertainty? An empirical evaluation of confidence elicitation in {LLMs}.
\newblock \emph{arXiv preprint arXiv:2306.13063}, 2023.

\bibitem[Zheng et~al.(2023)]{zheng2023judging}
L.~Zheng, W.-L. Chiang, Y.~Sheng, S.~Zhuang, Z.~Wu, Y.~Zhuang, Z.~Lin, Z.~Li, D.~Li, E.~P. Xing, H.~Zhang, J.~E. Gonzalez, and I.~Stoica.
\newblock Judging {LLM}-as-a-judge with {MT-Bench} and {Chatbot Arena}.
\newblock In \emph{NeurIPS}, 2023.

\bibitem[Zhou et~al.(2024)]{zhou2024webarena}
S.~Zhou, F.~F. Xu, H.~Zhu, X.~Zhou, R.~Lo, A.~Sridhar, X.~Cheng, T.~Ou, Y.~Bisk, D.~Fried, U.~Alon, and G.~Neubig.
\newblock {WebArena}: A realistic web environment for building autonomous agents.
\newblock In \emph{ICLR}, 2024.

\end{thebibliography}

\end{document}
